\section{Implementation}
\label{sec:impl}

We implement the pass as a single Java class, \texttt{com.jml.\allowbreak{}inferrer.\allowbreak{}analysis.\allowbreak{}CompositionalAnalyzer}, sitting next to the existing \texttt{InterproceduralAnalyzer} in the inferrer's module. The class is approximately 410 lines and depends only on the
inferrer's existing model (\texttt{MethodSpecification},
\texttt{SpecificationCache}, \texttt{CallGraph}) and JavaParser.

\subsection{Entry point}

The pass is invoked after the inferrer's isolated pass populates the
cache:

\begin{lstlisting}
SpecificationCache cache = new SpecificationCache();
CallGraph callGraph = new CallGraphBuilder()
        .buildFromCompilationUnits(units);
MethodSpecificationInferrer inferrer =
        new MethodSpecificationInferrer(cache, callGraph);
for (MethodDeclaration m : allMethods) {
    cache.put(signatureOf(m),
              inferrer.inferSpecification(m));
}

// Compositional refinement (this paper):
CompositionalAnalyzer.RefinementResult result =
    new CompositionalAnalyzer(cache, callGraph, declarations)
        .refineAll();
\end{lstlisting}

The \texttt{RefinementResult} carries the number of methods refined,
total iterations across all SCCs, and the SCC count. Callers can
choose whether to invoke the pass --- the design is opt-in so a
deployment that prefers the cheaper single-pass output can omit it.

\subsection{Two real bugs surfaced by the corpus-scale run}

We had validated the analyser with a 15-test unit suite before this
study. Running it on Apache Commons Lang surfaced two bugs that the
unit tests did not exercise.

\textbf{Bug 1: \texttt{lastIndexOf('.')} in \texttt{polymorphicCandidates}.}
The method searches for sibling cache entries by splitting each
signature into class name and method name. The split used
\texttt{sig.lastIndexOf('.')} --- which returns the position of the
\emph{rightmost} dot. For a signature like
\texttt{ArchUtils.getProcessor(java.lang.String)} the rightmost dot is
the one inside the parameter type \texttt{java.lang.String}, not the
separator between class and method. The resulting substring slice
threw \texttt{StringIndexOutOfBoundsException} with a negative-length
range. The fix is to use \texttt{lastIndexOf('.', paren - 1)}, which
restricts the search to before the parameter-list opening paren.

\textbf{Bug 2: same pattern in \texttt{resolveCallee}.} The same
algorithm and the same bug appeared in the call-resolution code.
Fixed identically.

Both bugs were silent under the unit-test inputs (which used short
parameter types) and only manifested on the corpus run with
fully-qualified parameter types from JavaParser. We report them
because they illustrate a general risk for analysers that work on
JavaParser ASTs: the printed representation of a type is
implementation-dependent, and any string-splitting code over
signatures must consider the worst case.

\subsection{Replication artefact}

The analyser, unit tests, per-library measurement harnesses, and
raw output logs are bundled in the anonymous replication package
linked from the title footnote. A single Maven invocation
(\texttt{mvn -Dtest=Compositional* test} on JDK~21) reproduces
every number in Section~\ref{sec:empirical} deterministically.
